\documentclass[11pt]{article}

\usepackage[margin=1in]{geometry}
\usepackage{amsmath,amssymb,amsthm,mathtools}
\usepackage{booktabs}
\usepackage{enumitem}
\usepackage{graphicx}
\usepackage{tikz}
\usetikzlibrary{arrows.meta,calc,positioning}
\usepackage[numbers,sort&compress]{natbib}
\usepackage{hyperref}
\usepackage{microtype}
\usepackage{xcolor}

\hypersetup{colorlinks=true,linkcolor=blue,citecolor=blue,urlcolor=blue}
\allowdisplaybreaks

\newtheorem{proposition}{Proposition}
\newtheorem{lemma}{Lemma}
\newtheorem{definition}{Definition}
\numberwithin{equation}{section}

\newcommand{\R}{\mathbb R}
\newcommand{\E}{\mathbb E}
\newcommand{\N}{\mathcal N}
\newcommand{\dd}{\,\mathrm d}
\newcommand{\tr}{\operatorname{tr}}
\newcommand{\diag}{\operatorname{diag}}
\newcommand{\vecop}{\operatorname{vec}}
\newcommand{\sgq}{\operatorname{SGQ}}
\newcommand{\chol}{\operatorname{chol}}
\newcommand{\ellhat}{\widehat\ell}
\newcommand{\bfone}{\mathbf 1}

\title{A Focus-Deepened Mathematical Companion To Fixed Sparse-Grid Quadrature Filtering And Its Analytical Gradient}
\author{Chak Wong}
\date{2026-06-12}

\begin{document}
\maketitle

\begin{abstract}
This note develops fixed sparse-grid quadrature filtering as a
mathematically explicit high-dimensional filtering proposal for the regime where
exact nonlinear filtering is unavailable but a deterministic Gaussian-surrogate
value-and-gradient target is still scientifically useful.  Starting from the
sparse-grid quadrature nonlinear filter of Jia, Xin, and Cheng
\cite{jia2012sparsegrid}, the note reconstructs the Gaussian approximation and
sparse-grid formulas in source order, then defines a fixed approximate
likelihood scalar together with an analytical derivative of that same scalar on
a declared branch.  The document is written for both review and implementation:
it states the exact filtering target, the narrower carried Gaussian surrogate,
the fixed cloud-construction rule, the per-step value recursion, the
same-scalar branch contract, the analytical gradient recursion, and the
finite-difference validity rule.  The central limitation is explicit.  The
reported scalar is not the exact nonlinear filtering likelihood, and the method
does not preserve general non-Gaussian posterior shape; it is a deterministic
Gaussian-surrogate lane whose usefulness depends on whether one carried Gaussian
and a fixed sparse-grid moment engine are adequate for the intended regime.
\end{abstract}

\section{Introduction And Scope}
\label{sec:p33-reader-orientation}

High-dimensional nonlinear filtering forces a choice of computational object.
The exact Bayesian filter propagates a full conditional law through time, but in
nonlinear models that law usually does not remain available in a finite closed
form.  In low dimension one may still tolerate heavy deterministic quadrature or
a richer non-Gaussian representation.  In high dimension, however, both cost and
structural complexity escalate quickly.  This note studies one narrower lane
suited to that regime: a fixed sparse-grid quadrature filter that carries a
single Gaussian surrogate, evaluates the required nonlinear moments by a fixed
deterministic cloud in standardized Gaussian coordinates, and differentiates the
same declared approximate scalar on a frozen computational branch.

The note is intended as a chapter-like companion to the monograph rather than as
a survey or a software manual.  Its task is to make one deterministic
Gaussian-surrogate lane precise enough to read, implement, audit, and compare.
The argument therefore moves from the exact filtering target, to the Gaussian
projection that replaces it, to the sparse-grid cloud that supplies the required
moments, and only then to the fixed scalar, its derivative, and the verification
obligations that make the lane auditable.

The central limitation should be fixed at the outset.  FixedSGQF does not claim
exact nonlinear filtering, exact nonlinear likelihood evaluation, or faithful
preservation of arbitrary multimodality, skewness, or high-order global
interaction structure.  Its purpose is narrower: to ask whether one carried
Gaussian, together with a fixed sparse-grid moment engine and a fixed-branch
value-and-gradient contract, can provide a coherent and scientifically useful
approximate likelihood lane in the regime under study.

\paragraph{Background assumed.}
The reader is assumed to know multivariate Gaussian notation, state-space-model
notation, Jacobian-style derivative notation, and the use of a Cholesky factor
for a positive-definite covariance matrix.  Sparse-grid merging, same-scalar
branch identities, and fixed-cloud derivative obligations are defined where they
first become necessary.

\tableofcontents

\subsection{Problem Setting And The FixedSGQF Lane}
\label{sec:p32-problem-lane}

The scientific problem is high-dimensional nonlinear filtering.  At time
\(t\), the exact Bayesian filter propagates and updates the full conditional law
\(p_\theta(x_t\mid y_{1:t})\).  In nonlinear models that law is generally not
available in a finite-dimensional closed form, and in high dimension it becomes
unrealistic to treat exact function-valued propagation as the practical
computational object.  FixedSGQF therefore does not attempt to rescue exactness
by notation.  It selects one deterministic approximation lane whose strengths and
boundaries can both be stated explicitly.

That lane is fixed sparse-grid quadrature filtering.  The carried state is a
single Gaussian surrogate \(\N(m_t,P_t)\).  The nonlinear moments needed by the
Gaussian projection are approximated with a deterministic sparse-grid rule in
standardized Gaussian coordinates.  The cloud, duplicate-merge policy,
covariance-factor branch, and veto rules are fixed before value and gradient
evaluation so that the reported derivative differentiates the same scalar that
the reported value routine computes.  The gain from that design is transparency:
each approximation stage is named, reproducible, and auditable.  The cost is
that a single Gaussian carried object cannot preserve general non-Gaussian
posterior geometry even when the quadrature rule is itself accurate for the
moments it is asked to approximate.

FixedSGQF should therefore not be read as a weaker version of a richer
density-approximation lane such as squared tensor-train filtering.  The present
note is about the deterministic Gaussian-surrogate lane with a fixed cloud and a
same-scalar analytical gradient.

\subsection{The Selected Approximation Lane In One Page}
\label{sec:p32-one-page}

The lane studied here is defined by three deliberate approximations.
\begin{enumerate}[leftmargin=0.28in]
  \item \textbf{Gaussian carried state.}  Instead of carrying the full filtering
  law, the method carries one Gaussian summary \(\N(m_t,P_t)\).
  \item \textbf{Sparse-grid moment closure.}  The nonlinear Gaussian moments that
  define the prediction and observation projection are approximated by a fixed
  sparse-grid rule in standardized coordinates.
  \item \textbf{Fixed-branch differentiation.}  The reported gradient is the
  derivative of a declared deterministic surrogate scalar on a frozen branch, not
  the derivative of a live adaptive algorithm.
\end{enumerate}

The resulting accumulated scalar is
\begin{equation}
  \ellhat_T^{\rm FSGQ}(\theta)
  =
  \sum_{t=1}^T \ellhat_t^{\rm FSGQ}(\theta),
  \label{eq:p32-one-page-scalar}
\end{equation}
where each increment is the Gaussian innovation log likelihood implied by the
sparse-grid moment approximation.  The exact object, the carried surrogate, and
the accumulated scalar are therefore different in kind:
\begin{equation}
\begin{aligned}
  \text{exact object:}
  &\quad p_\theta(x_t\mid y_{1:t}), \\
  \text{carried surrogate:}
  &\quad \N(m_t,P_t), \\
  \text{declared deterministic scalar:}
  &\quad \ellhat_T^{\rm FSGQ}(\theta).
\end{aligned}
  \label{eq:p32-exact-vs-approximate-summary}
\end{equation}
This distinction governs everything that follows.  The note does not return to it
repeatedly because it changes; it returns to it only when a later construction
needs the distinction in a sharper form.

\subsection{Approximation Hierarchy And Reading Guide}
\label{sec:p32-hierarchy}

The note unfolds in four large movements.  First, the exact filtering recursion
and the Gaussian projection that replaces it are fixed formally.  Second, the
sparse-grid construction is built from low-dimensional examples through
one-dimensional rules, tensor products, active index sets, and merged clouds.
Third, that fixed cloud is inserted into the filtering value recursion and then
into the same-scalar derivative recursion.  Fourth, the resulting lane is
verified, positioned against neighboring deterministic Gaussian alternatives, and
bounded by explicit non-claims.

The main coordinates are:
\begin{align}
  \xi &\in \R^{n_x},
  &\xi&\sim\N(0,I_{n_x})
  &&\hbox{standard quadrature coordinate},
  \label{eq:p32-standard-coordinate}\\
  x &= m+C\xi,
  &CC^\top&=P
  &&\hbox{physical state coordinate under a carried Gaussian},
  \label{eq:p32-physical-coordinate}\\
  a &= f_\theta(x),
  &&
  &&\hbox{transition image used for prediction},
  \label{eq:p32-transition-coordinate}\\
  z &= h_\theta(x),
  &&
  &&\hbox{observation image used for the innovation}.
  \label{eq:p32-observation-coordinate}
\end{align}
The sparse grid is built once in \(\xi\)-space and then placed into physical
state space through the current Gaussian factor.  That one sentence is worth
retaining before the technical development begins, because it is the bridge
between the cloud-construction chapter and the later value and gradient recursions.

\section{Exact Filtering And Gaussian Projection}
\label{sec:p31-ssm}

Use the additive-noise state-space model
\begin{align}
  X_t &= f_\theta(X_{t-1})+\eta_t,
  &
  \eta_t&\sim\N(0,Q_\theta),
  \label{eq:p31-state-model}\\
  Y_t &= h_\theta(X_t)+\varepsilon_t,
  &
  \varepsilon_t&\sim\N(0,R_\theta),
  \label{eq:p31-observation-model}
\end{align}
with \(X_t\in\R^{n_x}\), \(Y_t\in\R^{n_y}\), and parameter \(\theta\in\R^p\).
Throughout this note, the process noises \(\eta_t\) are assumed zero mean,
mutually independent of the observation noises \(\varepsilon_s\) for all relevant
time indices, independent across time within the process-noise family, and
independent of the past latent states and observations conditional on the model.
Likewise, the observation noises \(\varepsilon_t\) are assumed zero mean,
independent across time within the observation-noise family, and independent of
the latent states conditional on the declared additive observation model.  These
assumptions are exactly what let the filtering factorization and the later
covariance identities use \(Q_\theta\) and \(R_\theta\) as the declared additive
noise covariances.
The exact Bayesian prediction and update are
\begin{align}
  p_\theta(x_t\mid y_{1:t-1})
  &=
  \int p_\theta(x_t\mid x_{t-1})
       p_\theta(x_{t-1}\mid y_{1:t-1})\dd x_{t-1},
  \label{eq:p31-exact-prediction}\\
  p_\theta(x_t\mid y_{1:t})
  &=
  \frac{
  g_\theta(y_t\mid x_t)
  p_\theta(x_t\mid y_{1:t-1})}
  {\int g_\theta(y_t\mid x_t)
  p_\theta(x_t\mid y_{1:t-1})\dd x_t}.
  \label{eq:p31-exact-update}
\end{align}
These are the conceptual filtering equations used by Gaussian approximation
filters in Jia, Xin, and Cheng \cite[Section 2]{jia2012sparsegrid}.  In
nonlinear models, the integrals in \eqref{eq:p31-exact-prediction} and
\eqref{eq:p31-exact-update} usually do not close in finite-dimensional form.
Fixed SGQF replaces those exact densities by a Gaussian moment projection.

\subsection{Report-Wide Object Map And Implementation Conventions}
\label{subsec:p32-object-map}

The formal recursion below uses a stable set of objects throughout the value and
gradient paths.  The table is included here because these objects recur in later
chapters, but the implementation significance of each one is explained more
directly once the fixed cloud and branch record have been defined.
\begin{center}
\small
\begin{tabular}{p{0.20\linewidth}p{0.44\linewidth}p{0.22\linewidth}}
\toprule
symbol & meaning & shape \\
\midrule
\(\xi^{(r)}\) & standardized sparse-grid node in \(\R^{n_x}\) & \(n_x\) \\
\(w_r\) & accumulated signed weight after duplicate merge & scalar \\
\(m_t,P_t\) & carried Gaussian mean and covariance after update & \(n_x\), \(n_x\times n_x\) \\
\(m_t^-,P_t^-\) & predictive Gaussian mean and covariance before update & \(n_x\), \(n_x\times n_x\) \\
\(C_t, C_t^-\) & covariance factors with \(P_t=C_tC_t^\top\), \(P_t^-=C_t^-(C_t^-)^\top\) & \(n_x\times n_x\) \\
\(\chi_{t-1}^{(r)}\) & previous-state physical cloud point & \(n_x\) \\
\(a_t^{(r)}\) & transition image of \(\chi_{t-1}^{(r)}\) & \(n_x\) \\
\(\chi_t^{(r)}\) & predictive physical cloud point & \(n_x\) \\
\(z_t^{(r)}\) & observation image of \(\chi_t^{(r)}\) & \(n_y\) \\
\(\bar z_t\) & predicted observation mean & \(n_y\) \\
\(S_t\) & innovation covariance & \(n_y\times n_y\) \\
\(C_{xz,t}\) & state-observation cross-covariance & \(n_x\times n_y\) \\
\(v_t\) & innovation residual \(y_t-\bar z_t\) & \(n_y\) \\
\(K_t\) & Gaussian projection gain & \(n_x\times n_y\) \\
\(\ellhat_t^{\rm FSGQ}\) & one-step deterministic approximate log-likelihood increment & scalar \\
\bottomrule
\end{tabular}
\end{center}

Three implementation conventions are fixed report-wide.
\begin{enumerate}[leftmargin=0.28in]
  \item \textbf{Factor convention.}  Every covariance factor is the lower
  triangular Cholesky factor with positive diagonal, written
  \(C(P)=\chol(P)\), whenever the declared branch is valid.
  \item \textbf{Solve convention.}  Formulas may use \(S_t^{-1}\) for compact
  mathematics, but the implementation contract is to compute all such actions by
  linear solves against the declared Cholesky factor of \(S_t\), not by forming
  an explicit inverse.
  \item \textbf{Noise convention.}  Additive process noise enters analytically
  through the covariance term \(Q_\theta\) in predictive moments; additive
  observation noise enters analytically through \(R_\theta\) in the innovation
  covariance.  This report does not use state augmentation for these additive
  noises.
\end{enumerate}

When floating-point accumulation produces a covariance that is not exactly
symmetric, the report assumes explicit symmetrization by
\((P+P^\top)/2\) before branch checking.  If the resulting matrix is not
positive definite on the declared Cholesky branch, the branch is vetoed rather
than silently repaired.  Any jitter, clipping, or eigenvalue flooring policy
would define a different scalar than the one studied here.

\subsection{Gaussian Projection From Moments}
\label{sec:p31-gaussian-projection}

Assume the predictive state is represented by
\begin{equation}
  X_t^-\sim\N(m_t^-,P_t^-).
  \label{eq:p31-predictive-gaussian}
\end{equation}
Let
\begin{equation}
  Z_t=h_\theta(X_t^-)
  \label{eq:p31-predicted-observation}
\end{equation}
be the noiseless predicted observation.  The Gaussian projection needs
\begin{align}
  \bar z_t &= \E[Z_t],
  \label{eq:p31-zbar}\\
  S_t &= R_\theta+\E[(Z_t-\bar z_t)(Z_t-\bar z_t)^\top],
  \label{eq:p31-S}\\
  C_{xz,t}&=\E[(X_t^- -m_t^-)(Z_t-\bar z_t)^\top].
  \label{eq:p31-Cxz}
\end{align}
Given these moments, define
\begin{align}
  v_t&=y_t-\bar z_t,
  &
  K_t&=C_{xz,t}S_t^{-1},
  \label{eq:p31-v-K}\\
  m_t&=m_t^-+K_tv_t,
  &
  P_t&=P_t^- - K_tS_tK_t^\top.
  \label{eq:p31-projection-update}
\end{align}

\begin{proposition}[Affine projection]
\label{prop:p31-affine-projection}
Assume the innovation covariance \(S_t\) is square and invertible, so the solve
or inverse notation below is well defined, and all matrix products are conformable
under the stated object map.  Among estimators of \(X_t^-\) of the form
\(a+B Y_t\), the minimum mean-square estimator under the moments in
\eqref{eq:p31-zbar}--\eqref{eq:p31-Cxz} uses
\begin{equation}
  B=C_{xz,t}S_t^{-1},
  \qquad
  a=m_t^- - B\bar z_t.
  \label{eq:p41-affine-projection-intercept}
\end{equation}
Storing the resulting mean and covariance as a Gaussian gives
\eqref{eq:p31-projection-update}.
\end{proposition}

\begin{proof}
Center \(\widetilde X=X_t^- -m_t^-\) and
\(\widetilde Y=Y_t-\bar z_t\).  The squared-error objective for
\(m_t^-+B\widetilde Y\) is
\begin{equation}
  J(B)=\E\|\widetilde X-B\widetilde Y\|^2
  =
  \tr(P_t^-)-2\tr(BC_{xz,t}^\top)+\tr(BS_tB^\top).
  \label{eq:p31-affine-objective}
\end{equation}
Differentiating gives \(\nabla_BJ=-2C_{xz,t}+2BS_t\).  If \(S_t\) is
nonsingular, the stationary point is \(B=C_{xz,t}S_t^{-1}\).  Substitution gives
the residual covariance \(P_t^- - C_{xz,t}S_t^{-1}C_{xz,t}^\top\), which equals
\eqref{eq:p31-projection-update}.
\end{proof}

The approximation has two layers.  First, the moments
\eqref{eq:p31-zbar}--\eqref{eq:p31-Cxz} may be approximated by quadrature.
Second, the resulting posterior information is compressed to one Gaussian.

\subsection{How This Matches The Jia--Xin--Cheng Gaussian Approximation Block}
\label{subsec:p32-jia-gaussian-block}

Jia, Xin, and Cheng first write the Gaussian approximation filter before they
introduce sparse grids \cite[Section 2.2]{jia2012sparsegrid}.  In their
additive-noise setting, the prediction moment equations have the same
structure as
\begin{align}
  m_t^-&=\E[f_\theta(X_{t-1})],
  \label{eq:p32-source-pred-mean}\\
  P_t^-&=
  Q_\theta+
  \E[(f_\theta(X_{t-1})-m_t^-)(f_\theta(X_{t-1})-m_t^-)^\top],
  \label{eq:p32-source-pred-cov}
\end{align}
where \(X_{t-1}\sim\N(m_{t-1},P_{t-1})\) under the carried Gaussian
approximation.  The update moment equations have the structure
\begin{align}
  \bar z_t&=\E[h_\theta(X_t^-)],
  \label{eq:p32-source-zbar}\\
  S_t&=
  R_\theta+
  \E[(h_\theta(X_t^-)-\bar z_t)(h_\theta(X_t^-)-\bar z_t)^\top],
  \label{eq:p32-source-S}\\
  C_{xz,t}&=
  \E[(X_t^- -m_t^-)(h_\theta(X_t^-)-\bar z_t)^\top].
  \label{eq:p32-source-Cxz}
\end{align}
Then the Gaussian projection update is
\begin{align}
  K_t&=C_{xz,t}S_t^{-1},
  \label{eq:p32-source-K}\\
  m_t&=m_t^-+K_t(y_t-\bar z_t),
  \label{eq:p32-source-m-update}\\
  P_t&=P_t^- -K_tS_tK_t^\top.
  \label{eq:p32-source-P-update}
\end{align}
These equations are not yet sparse-grid equations.  They are the Gaussian moment
closure equations that any deterministic Gaussian quadrature lane must eventually
supply.  The next chapter now turns from that formal target to the concrete
low-dimensional construction of the sparse-grid cloud that will later be inserted
into the recursion.

\section{Low-Dimensional Construction Of FixedSGQF}
\label{sec:p34-motivating-cases}

The sparse-grid rule is easiest to understand before it is inserted into the full
filtering recursion.  The present chapter therefore begins with low-dimensional
constructions.  Its purpose is not to repeat the formal target, but to show what
kind of cloud is actually being built, why that cloud grows too quickly under a
naive tensor-product rule, and what the sparse-grid combination changes.

\subsection{One nonlinear Gaussian moment as the basic quadrature problem}
\label{subsec:p34-scalar-motivation}

Start with a scalar Gaussian input and a simple nonlinear observation map.  Let
\begin{equation}
  X\sim\N(m,P),
  \qquad
  h(X)=X^2.
  \label{eq:p34-scalar-gaussian-model}
\end{equation}
Suppose we want the Gaussian-projection observation mean
\begin{equation}
  \bar z = \E[h(X)] = \E[X^2].
  \label{eq:p34-scalar-target-mean}
\end{equation}
Write \(X=m+C\xi\) with \(\xi\sim\N(0,1)\) and \(C^2=P\).  Then
\begin{equation}
  \bar z
  =
  \int_{\R} (m+C\xi)^2\phi(\xi)\dd\xi.
  \label{eq:p34-standardized-integral}
\end{equation}
Once the state has been standardized, the moment problem becomes a
standard-normal weighted integral.  The quadrature rule is therefore not an
extra modeling choice layered on top of nowhere.  It is the numerical
approximation to the Gaussian expectation in
\eqref{eq:p34-standardized-integral}.

In this scalar example, a three-point symmetric rule already reproduces the
first few Gaussian moments exactly, which is why rules such as
\(\{-\sqrt3,0,\sqrt3\}\) appear naturally later.  The next step is to make that
one-dimensional rule explicit and then to see what happens when the same rule is
lifted into more than one coordinate.

\subsection{Tensor-product growth in two dimensions}
\label{subsec:p34-2d-motivation}

Now move from one coordinate to two.  Use a simple two-dimensional Gaussian
state,
\begin{equation}
  X\sim\N(m,P),
  \qquad
  X\in\R^2,
  \label{eq:p34-2d-gaussian}
\end{equation}
and imagine either a two-dimensional linear-Gaussian state-space step or a
slightly nonlinear observation moment built from that same Gaussian input.  The
point of this example is geometric rather than inferential: it shows how the
same standardized cloud is placed and why tensor products grow quickly.

Let \(X=m+C\xi\) with \(\xi\sim\N(0,I_2)\).  If each standardized coordinate
uses the three-point rule \(\{-\sqrt3,0,\sqrt3\}\), then the direct tensor
product uses all pairs
\begin{equation}
  (\xi_1,\xi_2)
  \in
  \{-\sqrt3,0,\sqrt3\}\times\{-\sqrt3,0,\sqrt3\},
  \label{eq:p34-2d-tensor-points}
\end{equation}
so already in two dimensions the rule has \(3^2=9\) points.  In three
dimensions the same construction has \(3^3=27\) points; in ten dimensions it has
\(3^{10}\) points.  The tensor-product rule is therefore transparent but already
expensive enough to motivate a more economical construction.

At this stage there is still no sparse-grid machinery.  Nothing has yet been
combined or pruned.  We are only taking one one-dimensional rule and using it in
each coordinate.  The tensor-product construction should therefore be read as the
rectangular standardized cloud that later sparse-grid ideas will modify, not as a
separate competing method.

\paragraph{What does the symbol \texorpdfstring{$F$}{F} mean?}
When the note later writes expressions such as \(I_{\mathbf i}[F]\) or
\(A_{b,L}[F]\), the symbol \(F\) is just the multivariate integrand being
averaged under the Gaussian weight.  In one dimension the same role is often
played by \(\psi\).  In several coordinates we simply switch to \(F(\xi)\) to
remind the reader that the function may depend on many variables.  A concrete
example to keep in mind is
\[
  F(\xi_1,\xi_2,\xi_3)=e^{\xi_1\xi_2\xi_3},
\]
which will be reused later to illustrate what a low sparse-grid level can and
cannot see.

\subsection{One-Dimensional Rules And Their Tensor Products}
\label{sec:p31-ghq}

The scalar example above shows why a Gaussian expectation becomes a
standard-normal weighted integral after standardization.  The next step is to fix
a concrete family of one-dimensional rules and then lift those rules into
multiple coordinates by tensor products.

\paragraph{What is a standard-normal Gauss--Hermite rule?}
A standard-normal Gauss--Hermite quadrature rule is a deterministic weighted-sum
approximation to a standard-normal integral,
\[
  \int_{\R}\psi(\xi)\phi(\xi)\dd\xi,
\]
chosen so that the rule is exact for as many low-degree polynomials as possible
for its number of points.  One picks nodes \(\zeta_{\ell,a}\) and weights
\(\omega_{\ell,a}\), evaluates the integrand at those nodes, and forms the
weighted sum
\[
  \sum_{a=1}^{s_\ell}\omega_{\ell,a}\psi(\zeta_{\ell,a}).
\]
Once a Gaussian moment has been standardized to the coordinate \(\xi\),
standard-normal GHQ gives a canonical deterministic way to approximate that
moment.

\paragraph{Which concrete family is being used in this note?}
At the teaching level, the easiest way to read the symbols is the following.  In
this note, the default BayesFilter family is the standard-normal GHQ ladder:
\[
  I_1=\hbox{1-point standard-normal GHQ},
  \qquad
  I_2=\hbox{3-point standard-normal GHQ},
  \qquad
  I_3=\hbox{5-point standard-normal GHQ},
  \qquad \ldots
\]
So the level label is a rule-family label, not a polynomial degree.  Level 2 does
not mean ``degree 2 polynomial.''  It means ``the second low-level member of the
chosen one-dimensional GHQ family,'' which in this note is the familiar symmetric
three-point rule.  The later formal policy statement \eqref{eq:p41-ghq-family}
records the same choice more precisely.

\paragraph{Where does the level-1 rule \texorpdfstring{$I_1:(0,1)$}{I1:(0,1)} come from?}
It is the simplest one-dimensional standard-normal rule: one node at the mean
\(\xi=0\) with total weight \(1\).  In the notation of the general one-dimensional
rule, this means \(s_1=1\), \(\zeta_{1,1}=0\), and \(\omega_{1,1}=1\), so
\[
  I_1[\psi]=\psi(0).
\]
This rule already integrates constants exactly, because
\[
  \int_{\R}1\cdot\phi(\xi)\dd\xi=1.
\]
So \(I_1:(0,1)\) is the natural level-1 base rule: evaluate at the center of
the standardized Gaussian and assign it the full mass.

\subsubsection{A fully worked one-dimensional toy rule}

Take the symmetric three-point ansatz
\begin{equation}
  \{(-a,\omega),(0,\omega_0),(a,\omega)\}.
  \label{eq:p36-three-point-ansatz}
\end{equation}
To match the first relevant Gaussian moments, require
\begin{align}
  \omega_0+2\omega &= 1,
  \label{eq:p36-three-point-mass}\\
  2\omega a^2 &= 1,
  \label{eq:p36-three-point-second}\\
  2\omega a^4 &= 3.
  \label{eq:p36-three-point-fourth}
\end{align}
Dividing \eqref{eq:p36-three-point-fourth} by
\eqref{eq:p36-three-point-second} gives \(a^2=3\).  Then
\(\omega=1/6\) and \(\omega_0=2/3\).  So the rule becomes
\begin{equation}
  \left\{\left(-\sqrt3,\frac16\right),\left(0,\frac23\right),
  \left(\sqrt3,\frac16\right)\right\}.
  \label{eq:p36-three-point-rule}
\end{equation}
This solved three-point rule is exactly the later level-2 rule \(I_2\) used in
the two-dimensional and three-dimensional examples.  In other words, \(I_2\) is
not a new assumption introduced later; it is simply shorthand for the symmetric
three-point standard-normal rule obtained from
\eqref{eq:p36-three-point-mass}--\eqref{eq:p36-three-point-fourth}.

\subsubsection{Lifting the rule to two dimensions}

The multi-index enters only because different coordinates may use different
one-dimensional rule levels.  Its role is bookkeeping, not new mathematics.  For
example:
\begin{itemize}[leftmargin=0.28in]
  \item \((1,1)\) means ``use level 1 in coordinate 1 and level 1 in coordinate
  2,''
  \item \((1,2)\) means ``use level 1 in coordinate 1 and level 2 in coordinate
  2,''
  \item \((2,1)\) means ``use level 2 in coordinate 1 and level 1 in coordinate
  2,''
  \item \((2,2)\) means ``use level 2 in both coordinates.''
\end{itemize}
Formally, for a multi-index \(\mathbf i=(i_1,\ldots,i_b)\), the corresponding
tensor rule is
\begin{equation}
  I_{\mathbf i}[F]
  =
  (I_{i_1}\otimes\cdots\otimes I_{i_b})[F].
  \label{eq:p31-tensor-rule}
\end{equation}
This means: apply the one-dimensional rule of level \(i_1\) in the first
coordinate, the one-dimensional rule of level \(i_2\) in the second coordinate,
and so on, then take the full Cartesian product of the resulting nodes and
weights.

In the concrete two-dimensional case, take \(b=2\) and use
\begin{equation}
  I_1:(0,1),
  \qquad
  I_2:\left\{\left(-\sqrt3,\frac16\right),\left(0,\frac23\right),
  \left(\sqrt3,\frac16\right)\right\}.
  \label{eq:p36-level1-level2-rules}
\end{equation}
Then the rule \((1,1)\) gives only \((0,0)\).  The rule \((1,2)\) gives the
three vertical points \((0,-\sqrt3)\), \((0,0)\), and \((0,\sqrt3)\).  The rule
\((2,1)\) gives the three horizontal points \((-\sqrt3,0)\), \((0,0)\), and
\((\sqrt3,0)\).  If both coordinates use level 2, then the full tensor product
rule \((2,2)\) gives the full nine-point grid
\begin{equation}
  (-\sqrt3,-\sqrt3),
  (-\sqrt3,0),
  (-\sqrt3,\sqrt3),
  (0,-\sqrt3),
  (0,0),
  (0,\sqrt3),
  (\sqrt3,-\sqrt3),
  (\sqrt3,0),
  (\sqrt3,\sqrt3),
  \label{eq:p35-nine-point-grid}
\end{equation}
with weights given by products of the one-dimensional weights.

The tensor-product rule is therefore simple but already expensive: if each
coordinate uses \(s\) nodes, then the point count is \(s^b\).  In two dimensions
with three points per coordinate, one already needs nine points; in three
dimensions the same direct construction would require twenty-seven.  The natural
next step is not to abandon the standardized-coordinate picture, but to reduce the
cost of that picture without losing the moment logic that made it useful.

\subsection{Three-dimensional preview: why sparse grids help}
\label{subsec:p37-3d-preview}

Use the same one-dimensional rules as before:
\begin{equation}
  I_1:(0,1),
  \qquad
  I_2:\left\{\left(-\sqrt3,\frac16\right),\left(0,\frac23\right),
  \left(\sqrt3,\frac16\right)\right\}.
  \label{eq:p37-level1-level2-rules}
\end{equation}
If all three coordinates use level 2, then the full tensor-product rule
\((2,2,2)\) places
\begin{equation}
  3\times 3\times 3 = 27
  \label{eq:p37-27-points}
\end{equation}
standardized points.  The sparse-grid idea begins by asking whether all those
point combinations are equally important at low level.

Now compare that with the sparse-grid band for \(b=3\) and \(L=2\).  This is a
different use of the word ``level.''  The tensor-product label \((2,2,2)\)
means ``use the one-dimensional level-2 rule in every coordinate.''  By
contrast, the sparse-grid level \(L=2\) means: combine all tensor rules whose
total index lies in the admissible band for \((b,L)=(3,2)\).  So the sparse-grid
construction does not keep the full 27-point tensor rule \((2,2,2)\).  Instead,
it mixes several lower-total-index tensor rules with signed coefficients.

The allowed multi-indices satisfy
\begin{equation}
  2\le |\mathbf i| \le 4.
  \label{eq:p37-3d-band}
\end{equation}
Since every coordinate level is at least 1, the only admissible three-dimensional
multi-indices are
\begin{equation}
  (1,1,1),
  \qquad
  (1,1,2),
  \qquad
  (1,2,1),
  \qquad
  (2,1,1).
  \label{eq:p37-3d-active-indices}
\end{equation}
So the sparse-grid rule does not use the full level-2 tensor product
\((2,2,2)\).  It instead combines one one-point tensor rule and three
one-axis three-point tensor rules.  This is the first low-dimensional picture of
why sparse grids help: they preserve the same standardized-coordinate logic while
paying first for lower-order interaction structure rather than for every corner
of the full tensor product.

\section{Filtering Value Path And Worked Example}
\label{sec:p31-value-path}

Fix a standardized sparse-grid cloud
\begin{equation}
  \mathcal C=\{(\xi^{(r)},w_r)\}_{r=1}^M,
  \qquad
  \xi^{(r)}\in\R^{n_x},
  \quad
  w_r\in\R.
  \label{eq:p31-cloud-for-filter}
\end{equation}
The cloud is frozen before likelihood evaluation.  The value path has two
quadrature stages at each time.

\paragraph{What this section gives the implementer.}
This section is the authoritative one-step value recursion.  A value routine
constructed from this note should take the current carried Gaussian,
current observation, fixed cloud, and declared model objects as input, then
return either an accepted one-step increment with updated \((m_t,P_t)\) or a
branch failure record.

\subsection{Prediction Stage}

Let \(P_{t-1}=C_{t-1}C_{t-1}^\top\) on the declared square-root branch.  Place
the previous Gaussian cloud:
\begin{equation}
  \chi_{t-1}^{(r)}=m_{t-1}+C_{t-1}\xi^{(r)}.
  \label{eq:p31-pred-placed-points}
\end{equation}
Propagate each point through the transition map,
\begin{equation}
  a_t^{(r)}=f_\theta(\chi_{t-1}^{(r)}),
  \label{eq:p31-transition-points}
\end{equation}
and form the predictive moments
\begin{align}
  m_t^- &= \sum_{r=1}^{M} w_r a_t^{(r)},
  \label{eq:p31-pred-mean}\\
  P_t^- &= Q_\theta +
  \sum_{r=1}^{M}w_r(a_t^{(r)}-m_t^-)(a_t^{(r)}-m_t^-)^\top.
  \label{eq:p31-pred-cov}
\end{align}
If the symmetrized \(P_t^-\) fails the declared Cholesky branch, the value path
returns a predictive-branch failure record rather than a scalar.

\subsection{Observation Stage}

Assume the predictive branch is valid, so that
\begin{equation}
  P_t^-=C_t^-(C_t^-)^\top.
  \label{eq:p31-pred-factor}
\end{equation}
Place the predictive cloud,
\begin{equation}
  \chi_t^{(r)}=m_t^-+C_t^-\xi^{(r)},
  \label{eq:p31-obs-placed-points}
\end{equation}
propagate it through the observation map,
\begin{equation}
  z_t^{(r)}=h_\theta(\chi_t^{(r)}),
  \label{eq:p31-obs-values}
\end{equation}
and compute the observation moments
\begin{align}
  \bar z_t &= \sum_{r=1}^{M}w_r z_t^{(r)},
  \label{eq:p31-obs-mean}\\
  S_t &= R_\theta+
  \sum_{r=1}^{M}w_r(z_t^{(r)}-\bar z_t)(z_t^{(r)}-\bar z_t)^\top,
  \label{eq:p31-S-cloud}\\
  C_{xz,t} &= \sum_{r=1}^{M}w_r(\chi_t^{(r)}-m_t^-)(z_t^{(r)}-\bar z_t)^\top.
  \label{eq:p31-cross-cov}
\end{align}
Then define
\begin{align}
  v_t &= y_t-\bar z_t,
  \label{eq:p31-innovation}\\
  K_t &= C_{xz,t}S_t^{-1},
  \label{eq:p31-kalman-gain}\\
  m_t &= m_t^-+K_t v_t,
  \label{eq:p31-filter-mean}\\
  P_t &= P_t^- - K_tS_tK_t^\top.
  \label{eq:p31-filter-cov}
\end{align}
The one-step deterministic approximate log-likelihood increment is
\begin{equation}
  \ellhat_t^{\rm FSGQ}
  =
  -\frac12
  \left\{
    \log\det S_t + v_t^\top S_t^{-1}v_t + n_y\log(2\pi)
  \right\}.
  \label{eq:p31-fsgq-loglik}
\end{equation}
The accumulated scalar is the sum of these increments over time.

\subsection{Worked One-Step Example With A Numeric Oracle}
\label{sec:p32-worked-example}

Use the scalar nonlinear model
\begin{equation}
  X_1 = \beta X_0 + \eta_1,
  \qquad
  \eta_1\sim\N(0,Q),
  \qquad
  Y_1 = X_1^2 + \varepsilon_1,
  \qquad
  \varepsilon_1\sim\N(0,R),
  \label{eq:p32-example-model}
\end{equation}
with
\begin{equation}
  X_0\sim\N(m_0,P_0),
  \qquad
  m_0=0.5,
  \qquad
  P_0=1,
  \qquad
  \beta=1,
  \qquad
  Q=4,
  \qquad
  R=1,
  \qquad
  y_1=3.
  \label{eq:p32-example-params}
\end{equation}
Use the three-point level-2 cloud
\[
  \xi\in\{-\sqrt3,0,\sqrt3\},
  \qquad
  w\in\{1/6,2/3,1/6\}.
\]
Then
\begin{align}
  m_1^- &= \beta m_0 = 0.5,
  &P_1^- &= \beta^2P_0+Q = 5,\\
  C_1^- &= \sqrt5,
  &\chi_1^{(r)} &\in \{0.5-\sqrt{15},\ 0.5,\ 0.5+\sqrt{15}\}.
\end{align}
The observation values are
\begin{equation}
  z_1^{(r)} = (\chi_1^{(r)})^2
  \in
  \left\{\frac{31}{2}-\sqrt{15},\ \frac14,\ \frac{31}{2}+\sqrt{15}\right\}.
  \label{eq:p32-example-z-values}
\end{equation}
From these,
\begin{align}
  \bar z_1&=2.5=\frac52,
  &S_1&=5.375=\frac{43}{8},\\
  C_{xz,1}&=1.25=\frac54,
  &v_1&=0.5=\frac12,\\
  K_1&=\frac{10}{43}\approx 0.2325581395,
  &m_1&=\frac{53}{86}\approx0.6162790698,\\
  P_1&=\frac{165}{172}\approx0.9593023256.
  \label{eq:p32-example-moments}
\end{align}
The one-step deterministic approximate log likelihood is
\begin{equation}
  \ellhat_1^{\rm FSGQ}
  =
  -\frac12\left\{\log S_1+v_1^2S_1^{-1}+\log(2\pi)\right\}
  \approx -1.7830736342.
  \label{eq:p32-example-ell}
\end{equation}
This is the value-path oracle used later for both implementation checking and
same-branch derivative verification.

\section{Same-Scalar Branch Contract And Analytical Gradient}
\label{sec:p31-same-scalar}

The branch record \(\mathcal B\) freezes both the value-defining structural
choices and the derivative-support metadata needed to differentiate the same
scalar on that branch.  The value scalar itself depends only on the former, while
the derivative-support objects matter only when the gradient lane is requested.
A fixed-SGQF branch is the tuple
\begin{equation}
\begin{aligned}
  \mathcal B
  =
  (&y_{1:T},\ \text{observation preprocessing},\
  m_0(\theta),P_0(\theta),\dot m_0,\dot P_0,\\
  &\mathcal C,\ \text{one-dimensional rules},\
  \text{duplicate-merge tolerance},\
  \text{zero-weight rule},\
  \text{node ordering},\\
  &C(P)=\chol(P),\
  \text{symmetrize-then-veto rule},\
  \epsilon_S,\epsilon_P,\\
  &f_\theta,h_\theta,Q_\theta,R_\theta,\
  D_xf_\theta,\partial_i f_\theta,
  D_xh_\theta,\partial_i h_\theta,\dot Q_\theta,\dot R_\theta).
\end{aligned}
  \label{eq:p31-branch-tuple}
\end{equation}
The fixed scalar is
\begin{equation}
  \ellhat_T^{\rm FSGQ}(\theta;\mathcal B)
  =
  \sum_{t=1}^T \ellhat_t^{\rm FSGQ}(\theta;\mathcal B),
  \label{eq:p31-fixed-scalar}
\end{equation}
where every term is computed by
\eqref{eq:p31-pred-placed-points}--\eqref{eq:p31-filter-cov}.  Changing the
grid level, merge tolerance, node ordering, factor branch, or veto rule changes
the scalar.  The derivative formulas below therefore cover the open branch where
the declared branch remains valid without repair.

\paragraph{Plain-language branch rule.}
The branch identity freezes structure, not numerical values.  When \(\theta\)
changes inside a same-scalar finite-difference check, the moments, gains, and
likelihood contributions are recomputed, but they must be recomputed by the same
cloud, the same merge/order policy, the same factor family, and the same realized
acceptance path through the branch logic.  This is the operational meaning of
``same scalar'' for implementation review.

\subsection{Analytical Gradient Of The Fixed Scalar}
\label{sec:p31-gradient}

The derivative chapter follows the same order as the value chapter.  Differentiate
cloud placement, transition images, predictive moments, observation images,
innovation statistics, and finally the innovation log-likelihood increment, all
while holding the branch record fixed.

\paragraph{Why this derivative is operationally delicate.}
The scalar is deterministic only on a declared branch.  The derivative is
therefore not merely the formal derivative of some symbolic quadrature formula.
It is the derivative of the declared deterministic scalar evaluated on the same
branch.  If a nearby perturbation changes the merge rule, factor branch, or veto
path, the finite difference has left the scalar being differentiated.

\subsubsection{The Story Of The Derivative}
The derivative propagation is easiest to read in the same order as the value path:
place the cloud, push it through the transition, form predictive moments, place
again, push through the observation, form innovation statistics, then
accumulate the score.  The detailed formulas and sensitivity identities from the
current report remain unchanged and are retained below as the formal derivative
contract.

\subsubsection{Square-Root Branch}
\label{subsec:p31-square-root-branch}

The declared covariance factor is the lower-triangular Cholesky factor.  Its
matrix differential convention, derivative identities, and subsequent sensitivity
propagation rules are the same as in p44 and remain the authoritative branch
specification here.

\subsubsection{Prediction Sensitivities}
\label{subsec:p31-prediction-sensitivities}

The prediction sensitivities are obtained by differentiating the cloud placement,
transition images, predictive mean, and predictive covariance under the fixed
cloud and branch record.  The formulas retained from the current report remain
unchanged and continue to define the same-scalar derivative contract.

\subsubsection{Observation Sensitivities}
\label{subsec:p31-observation-sensitivities}

The observation sensitivities are obtained by differentiating the predictive cloud
placement, observation images, predicted observation mean, innovation covariance,
and cross-covariance under the same branch.

\subsubsection{Innovation Score}
\label{subsec:p31-innovation-score}

The innovation score is the derivative of the one-step Gaussian innovation log
likelihood under the declared branch.  The detailed formulas from the current
report remain the operational score definition.

\subsubsection{Posterior Sensitivity Propagation}
\label{subsec:p31-posterior-sensitivity}

Once the score, gain, and covariance sensitivities have been computed, propagate
the posterior mean and covariance sensitivities in the same order as the value
path.  The report's existing derivative formulas remain the authoritative
implementation contract.

\subsection{One Boxed Mathematical Algorithm}
\label{sec:p31-boxed-algorithm}

The boxed algorithm remains the chapter summary of the full same-scalar value and
gradient recursion.  In p45 it remains at the end of the derivative chapter so
that it reads as a culmination rather than as a competing chapter line.

\section{Verification And Validation}
\label{sec:p31-fd}

The fixed-SGQF lane is only useful if value, branch, and derivative claims can be
verified on the same scalar.  This chapter therefore gathers the finite-difference
check, the worked toy traces, and the diagnostic test models into one verification
ledger.  The central question here is not merely whether the code runs, but
whether the mathematical object declared earlier is the same object being tested,
differentiated, and compared.

\subsection{Finite-Difference Same-Scalar Check}
\label{subsec:p31-finite-difference}

The same-scalar finite-difference check asks one narrow question: if the branch
identity remains unchanged under a small parameter perturbation, does the
analytical gradient agree with the finite difference of the same declared scalar?
The detailed formulas and thresholds from the current report remain unchanged; the
point of the present reorganization is only to keep them in one verification
chapter.

\paragraph{Implementation-facing FD pseudocode summary.}
The finite-difference routine should perturb parameters only inside the same
branch, re-evaluate the fixed scalar on that same branch identity, and reject any
comparison that silently changes the structural scalar.

\subsection{Diagnostics, Accuracy, Memory, And Performance Tests}
\label{sec:p31-validation}

The diagnostic chapter keeps the same hierarchy of checks as before, but now reads
as one verification ledger rather than as an independent companion note.  Use the
linear-Gaussian recovery case to check exact collapse, the scalar quadratic case
to check low-order moment correctness, the cloud-sensitive nonlinear case to check
the effect of cloud placement, the block-local high-dimensional case to test the
intended sparse-grid advantage, and the hard all-coordinate interaction case to
mark the lane's scientific boundary.

\paragraph{Model A: linear-Gaussian recovery.}
This is the exact-collapse test.  FixedSGQF should match the Kalman recursion when
the model is linear-Gaussian.

\paragraph{Model B: scalar quadratic observation.}
Use \eqref{eq:p31-quadratic-cell} as the smallest nonlinear moment test that still
admits a transparent oracle.
\begin{equation}
  X\sim\N(0,P),\qquad Y=X^2+\varepsilon,\qquad
  \varepsilon\sim\N(0,R).
  \label{eq:p31-quadratic-cell}
\end{equation}

\paragraph{Model C: cloud-sensitive two-dimensional nonlinear observation.}
This is the first test where cloud geometry, not only algebraic consistency,
matters.

\paragraph{Model D: block-local high-dimensional nonlinear model.}
This is the first intended operating regime for a moderate sparse-grid level.

\paragraph{Model E: deliberately hard all-coordinate interaction.}
This is the regime in which the Gaussian-surrogate lane becomes too narrow and the
note should fail honestly rather than overclaim.

\paragraph{Model F: Jia--Xin--Cheng orbit-style benchmark.}
This preserves the source-linked large-scale benchmark lane already present in the
report.

\section{Adaptive Sparse Grids As Grid Design}
\label{sec:p31-adaptive-grid-design}

Adaptive sparse grids remain relevant in this report only as an offline grid-design
remark.  The present fixed-scalar derivative contract does not differentiate a
live adaptive selection rule.  So adaptation belongs here as a late remark on grid
design, not as a competing main-line method.

\section{Relation To Neighboring High-Dimensional Filtering Proposals}
\label{sec:p31-comparison}

The present lane should now be positioned only after its object, cloud,
recursion, derivative, and verification obligations are already clear.  The main
comparison question is not whether FixedSGQF dominates every neighboring method.
It is where this deterministic Gaussian-surrogate lane sits among other serious
high-dimensional proposals and what exact trade it makes to earn its place.

The UKF comparison should now be read mainly by cross-reference back to the
low-dimensional construction chapter.  The key message has already been stated
there: low-level SGQF can look UKF-like geometrically and may even share final
weights in special symmetric-axis cases, but the sparse-grid construction is
assembled from one-dimensional rules, active tensor-rule selection, and Smolyak
recombination rather than from one directly specified sigma-point formula.

The remaining positioning arguments against EKF, dense GHQF, live adaptive sparse
grids, particle methods, and density-based tensor methods are preserved from p44,
but now read as late synthesis rather than as repeated restarts of the main lane
claim.

\section{Conclusion}
\label{sec:p31-conclusion}

FixedSGQF is a deliberately narrow high-dimensional filtering lane.  It does not
attempt to preserve a full non-Gaussian posterior law.  It carries one Gaussian
surrogate, computes the required nonlinear moments by a fixed sparse-grid cloud,
and differentiates the same declared deterministic scalar on a frozen branch.
That narrowness is the point.  It is what makes the method reproducible,
auditable, and mathematically explicit enough to support implementation and
verification.

The note therefore closes with a bounded claim.  It does not argue that FixedSGQF
is the best nonlinear filter overall.  It argues that, when a deterministic
Gaussian-surrogate lane with a same-scalar analytical gradient is the desired
object, the fixed sparse-grid construction gives a coherent and technically
well-specified way to build that lane.

\appendix

\section{Implementation Contract}
\label{sec:p31-implementation-contract}

This chapter records the implementation-facing material only after the reader has
seen the formal target, the cloud construction, the value recursion, and the
branch-conditioned derivative.  The point is not merely to list contracts.  It is
to state exactly what the runtime object is, which decisions change the scalar,
and which conventions an implementation can silently get wrong.  The detailed
input-output contracts, failure conventions, and runtime invariants remain
unchanged from the earlier notes; what changes here is the emphasis.  The reader
should now know why these details matter.

\section{End-To-End Mathematical Algorithm}
\label{sec:p32-end-to-end-algorithm}

The end-to-end algorithm remains as appendix material so that the main text can
carry the mathematical argument while the appendix records the full execution
ordering explicitly.  Read it here not as a second explanation of the method, but
as the exact ordering that preserves the cloud identity, branch identity, and
same-scalar value-and-gradient contract established earlier.

\section{Where Each Construction Comes From}
\label{sec:p32-source-map}

The source map remains useful as an audit appendix linking the BayesFilter
construction back to the underlying literature and internal specialization
choices.

\section{Notation And Formula Inventory}
\label{sec:p32-defaults}

The notation and formula inventory remains as late technical support rather than
as a main-text interruption.

\bibliographystyle{plainnat}
\bibliography{../references}

\end{document}
